% =============================================================
% Chapitre I — Fondements du downscaling climatique
% Style cible : gabarit Donald FEUZING NTEMMA (URIFIA 2024)
% =============================================================

\documentclass[14pt,a4paper,openany]{extreport}

% ---------- Encodage et langue ----------
\usepackage[T1]{fontenc}
\usepackage[utf8]{inputenc}
\usepackage[french]{babel}

% ---------- Police lisible pour impression papier ----------
\usepackage{lmodern}

% ---------- Interligne (impression papier reliée) ----------
\usepackage{setspace}
\onehalfspacing

% ---------- Géométrie (marges adaptées à la reliure) ----------
\usepackage[a4paper,inner=3cm,outer=2.3cm,top=2.5cm,bottom=2.7cm,headheight=15pt,bindingoffset=0.3cm]{geometry}

% ---------- Mathématiques ----------
\usepackage{amsmath,amssymb,amsthm}

% ---------- Graphiques et couleurs ----------
\usepackage{graphicx}
\usepackage{xcolor}
\definecolor{uridarkblue}{RGB}{30,40,90}
\definecolor{urigray}{RGB}{70,70,70}

% ---------- Boîte SOMMAIRE et placeholders de figures ----------
\usepackage[most]{tcolorbox}

% ---------- Mise en forme des chapitres et sections ----------
% Numérotation : chiffres romains pour les chapitres, I.X.Y pour le reste
\renewcommand{\thechapter}{\Roman{chapter}}
\renewcommand{\thesection}{\thechapter.\arabic{section}}
\renewcommand{\thesubsection}{\thesection.\arabic{subsection}}
\renewcommand{\thesubsubsection}{\thesubsection.\arabic{subsubsection}}

% Redéfinition manuelle de \@makechapterhead pour le style "Chapitre I + titre"
% (plus robuste que titlesec quand on a déjà tcolorbox+amsmath+hyperref).
\makeatletter
\renewcommand{\@makechapterhead}[1]{%
  \vspace*{20\p@}%
  {\parindent \z@ \centering \normalfont
    \textsc{\large Chapitre}\par
    \vspace{0.3em}
    {\color{uridarkblue}\fontsize{60}{72}\selectfont\bfseries\thechapter\par}
    \vspace{1em}
    {\color{uridarkblue}\Large\bfseries\MakeUppercase{#1}\par}
    \vspace{1em}
    \rule{0.4\linewidth}{0.6pt}\par
    \vspace{2em}
  }}
\renewcommand{\@makeschapterhead}[1]{%
  \vspace*{20\p@}%
  {\parindent \z@ \centering \normalfont
    {\color{uridarkblue}\Large\bfseries\MakeUppercase{#1}\par}
    \vspace{1em}
    \rule{0.4\linewidth}{0.6pt}\par
    \vspace{2em}
  }}
\makeatother

% Sections, sous-sections, sous-sous-sections via titlesec
\usepackage{titlesec}
\titleformat{\section}
  {\normalfont\Large\bfseries\color{uridarkblue}}
  {\thesection.}{0.8em}{}
\titleformat{\subsection}
  {\normalfont\large\bfseries}
  {\thesubsection.}{0.7em}{}
\titleformat{\subsubsection}
  {\normalfont\normalsize\bfseries\itshape}
  {\thesubsubsection.}{0.6em}{}

% ---------- En-têtes et pieds de page ----------
\usepackage{fancyhdr}
\pagestyle{fancy}
\fancyhf{}
\fancyhead[L]{\textsc{\nouppercase{\leftmark}}}
\fancyhead[R]{\thepage}
\fancyfoot[L]{\footnotesize Mémoire de Master en Informatique, Université de Dschang}
\fancyfoot[R]{\footnotesize URIFIA}
\renewcommand{\headrulewidth}{0.4pt}
\renewcommand{\footrulewidth}{0.2pt}
\renewcommand{\chaptermark}[1]{\markboth{#1}{}}
\renewcommand{\sectionmark}[1]{\markright{\thesection\ #1}{}}

% ---------- Légendes de figures / tableaux ----------
\usepackage[font=small,labelfont={sc,bf},textfont=it]{caption}
\renewcommand{\thefigure}{\arabic{figure}}
\renewcommand{\thetable}{\arabic{table}}

% ---------- Citations numériques ----------
\usepackage[numbers,square,sort&compress]{natbib}

% ---------- Hyperref (en dernier) ----------
\usepackage[hidelinks]{hyperref}

% ---------- Macros pratiques ----------
\newcommand{\Oracle}{\textsc{Oracle}}
\newcommand{\CorrDiff}{\textsc{CorrDiff}}
\newcommand{\LR}{\ensuremath{\text{LR}}}
\newcommand{\HR}{\ensuremath{\text{HR}}}
\newcommand{\xHR}{\ensuremath{x_{\HR}}}
\newcommand{\yLR}{\ensuremath{y_{\LR}}}
\newcommand{\muHR}{\ensuremath{\mu_{\HR}}}

% Boîte SOMMAIRE personnalisée
\newtcolorbox{sommairebox}{
  enhanced,
  colback=white,
  colframe=urigray,
  boxrule=0.5pt,
  arc=0pt,
  left=1.5em,right=1.5em,top=0.8em,bottom=0.8em,
  title={\textsc{Sommaire}},
  attach boxed title to top center={yshift=-2mm,xshift=0mm},
  boxed title style={
    colback=white,colframe=white,coltitle=urigray,boxrule=0pt,
    left=0.8em,right=0.8em,top=0pt,bottom=0pt
  },
  fonttitle=\bfseries\small,
}

% Placeholder de figure (utilisé tant que les figures ne sont pas produites)
\newcommand{\figplaceholder}[2][0.9\linewidth]{%
  \begin{tcolorbox}[width=#1,height=5cm,valign=center,colback=gray!8,colframe=gray!45,arc=2pt,boxrule=0.5pt]
    \centering\small\textit{#2}
  \end{tcolorbox}%
}

% Labels de chapitres futurs (pour éviter les ?? dans les renvois)
\providecommand{\chaplabel}[1]{}

\setcounter{page}{6}

\begin{document}

\chapter{Fondements du downscaling climatique}
\label{chap:fondements}

% ---------- Boîte SOMMAIRE manuelle ----------
\begin{sommairebox}
\small
\noindent I.1 \ Introduction \dotfill \pageref{sec:intro}\\
\noindent I.2 \ Contexte climatique et besoin de désagrégation \dotfill \pageref{sec:contexte}\\
\noindent I.3 \ Données climatiques et leurs propriétés \dotfill \pageref{sec:donnees}\\
\noindent I.4 \ Approches classiques de downscaling \dotfill \pageref{sec:classiques}\\
\noindent I.5 \ Approches probabilistes et génératives \dotfill \pageref{sec:generatifs}\\
\noindent I.6 \ Causalité et modèles structurels \dotfill \pageref{sec:causalite}\\
\noindent I.7 \ Verrous scientifiques et trilemme \dotfill \pageref{sec:verrous}\\
\noindent I.8 \ Synthèse \dotfill \pageref{sec:synthese}
\end{sommairebox}

\vspace{1.5em}

% =============================================================
\section{Introduction}
\label{sec:intro}
% =============================================================

L'intensification du changement climatique et la fréquence croissante d'événements extrêmes imposent aux communautés locales --- en particulier en Afrique centrale et au Cameroun --- des décisions de plus en plus fines en matière d'agriculture, de gestion de l'eau et de planification urbaine~\cite{ipcc_ar6_2021}. Ces décisions requièrent des projections climatiques à haute résolution spatiale, idéalement au kilomètre près. Or les modèles globaux disponibles (CMIP6~\cite{eyring_cmip6_2016}) opèrent à la centaine de kilomètres. Combler ce fossé est l'objet du \emph{downscaling} climatique, dont une variante moderne, fondée sur l'apprentissage génératif, fait l'objet du présent mémoire.

Ce chapitre établit les fondements nécessaires à la compréhension de notre contribution : contexte climatique et besoin de désagrégation (section~\ref{sec:contexte}), données mobilisées (section~\ref{sec:donnees}), méthodes classiques (section~\ref{sec:classiques}), approches probabilistes et génératives, notamment le paradigme à deux étages \CorrDiff{}~\cite{mardani_corrdiff_2023} (section~\ref{sec:generatifs}), causalité structurelle pour \Oracle{} (section~\ref{sec:causalite}), verrous scientifiques (section~\ref{sec:verrous}) et synthèse (section~\ref{sec:synthese}).

% =============================================================
\section{Contexte climatique et besoin de désagrégation}
\label{sec:contexte}
% =============================================================

\subsection{Changement climatique et impacts régionaux}
\label{ssec:cc-regional}

Le sixième rapport du GIEC (2021) établit une hausse d'environ $1{,}1\,^{\circ}\mathrm{C}$ de la température moyenne globale par rapport à 1850--1900, et une accélération de la tendance depuis 1980~\cite{ipcc_ar6_2021}. Au-delà de la moyenne, ce sont surtout l'intensification des pluies torrentielles, l'allongement des sécheresses et le bouleversement des régimes de mousson qui retiennent l'attention.

Ces transformations sont inégalement réparties. En Afrique centrale et de l'Ouest, les projections font état d'extrêmes pluvieux accrus et d'une variabilité interannuelle renforcée~\cite{ipcc_ar6_2021,almazroui_2020}, avec des conséquences déjà mesurables sur les rendements agricoles, les cycles hydrologiques et les risques urbains d'inondation. Anticiper ces évolutions à l'échelle locale --- semis adaptés en zone sahélienne, dimensionnement de l'assainissement à Yaoundé ou Douala --- est précisément la motivation de la désagrégation spatiale.

\subsection{Modèles climatiques globaux et régionaux}
\label{ssec:gcm-rcm}

Un \textbf{modèle climatique global} (\emph{General Circulation Model}, GCM) résout numériquement les équations de la dynamique atmosphérique et océanique pour produire des projections multi-décennales sous divers scénarios d'émissions. Le projet CMIP6 coordonne plus d'une centaine de tels modèles, dont ACCESS-CM2 (CSIRO, Australie)~\cite{eyring_cmip6_2016}, à des résolutions de $100$ à $250\,\mathrm{km}$ limitées par le coût numérique.

À cette échelle, les processus locaux (convection profonde, effets orographiques, brises côtières, mésoéchelle) sont \emph{paramétrés} plutôt qu'explicitement résolus, ce qui introduit des biais structurels marqués pour la précipitation. Les \textbf{modèles climatiques régionaux} (RCM) résolvent les mêmes équations à $10$--$50\,\mathrm{km}$ sur un domaine restreint, emboîtés dans un GCM aux frontières : ce \emph{downscaling dynamique} améliore le réalisme régional mais reste très coûteux (plusieurs millions d'heures de calcul par scénario). Le passage de la grille GCM ($\sim$$250\,\mathrm{km}$) à la grille cible HR ($\sim$$5\,\mathrm{km}$) représente un facteur d'environ 50 par dimension.
\label{fig:resolutions}

\subsection{Besoins locaux et cas de la précipitation}
\label{ssec:besoin-local}

L'écart entre la résolution disponible et celle requise par les utilisateurs --- le \emph{gap de résolution} --- est critique en agriculture (calendriers culturaux, irrigation), en hydrologie (forçages pluie-débit pour la prévision des crues) et en planification urbaine (assainissement, cartographie des zones inondables), tous trois exigeant des champs au kilomètre. Le contexte africain est particulièrement contraint : CORDEX-Africa fournit des RCM à $25$--$50\,\mathrm{km}$ mais reste lacunaire et hors de portée des équipes locales pour des scénarios multi-modèles plus fins.

Ce constat motive le \emph{downscaling statistique profond}, qui apprend, à partir d'un corpus historique de couples LR/HR, une fonction probabiliste $\xHR \sim p(\cdot\,|\,\yLR)$ déployable à coût marginal une fois entraînée.
\label{ssec:cas-precipitation}
Parmi les variables candidates, la \textbf{précipitation} occupe une place centrale : c'est la variable la plus liée aux risques perçus par les populations (sécheresse, inondations, sécurité alimentaire) et la plus exigeante statistiquement --- forte intermittence, queue lourde, hétérogénéité spatiale~\cite{maraun_widmann_2018,vandal_deepsd_2017} ---, si bien qu'une méthode qui y réussit se transfère généralement aux variables plus régulières (température, pression). Le présent mémoire la retient comme variable d'illustration et de validation. \Oracle{} est entraîné et évalué sur des précipitations quotidiennes du couple ACCESS-CM2 / Nouvelle-Zélande, terrain expérimental qui combine cible HR de qualité et baseline génératif de référence dans la littérature. L'architecture est conçue pour être étendue à un cadre multivariable, son encodeur traitant déjà plusieurs variables de conditionnement ; cette extension est laissée aux travaux futurs.

% =============================================================
\section{Données climatiques et leurs propriétés}
\label{sec:donnees}
% =============================================================

\subsection{Sources de données}
\label{ssec:sources}

Trois catégories alimentent le downscaling. Les \textbf{observations directes} (stations météorologiques) offrent une grande qualité ponctuelle mais une couverture spatiale très inégale, particulièrement lacunaire en Afrique et dans les zones océaniques~\cite{wmo_gcos_2022}. Les \textbf{réanalyses atmosphériques} pallient cette lacune en assimilant en continu observations stations, radiosondages et satellites dans un modèle numérique : ERA5~\cite{hersbach_era5_2020} fournit ainsi des champs horaires à $\sim$$31\,\mathrm{km}$ depuis 1940 et constitue la référence gridded historique. Les \textbf{simulations climatiques} (GCM et RCM, sorties CMIP6~\cite{eyring_cmip6_2016}) couvrent la période historique (1850--2014) et des projections futures sous scénarios SSP. Dans notre travail, ACCESS-CM2 fournit le conditionnement basse résolution et un produit régional néo-zélandais la cible HR.

\subsection{Variables, formats et conventions}
\label{ssec:formats}

Les données climatiques sont distribuées au format \texttt{NetCDF} (conventions \emph{Climate and Forecast}), auto-descriptif (tableaux + métadonnées). Conformément à CMIP6, la précipitation totale \texttt{pr} est exprimée en $\mathrm{kg}\,\mathrm{m}^{-2}\,\mathrm{s}^{-1}$ et convertie en mm/jour pour l'analyse. Le conditionnement atmosphérique mobilise \texttt{tas}, \texttt{tasmin}, \texttt{tasmax} (températures à 2~m), \texttt{sfcWind} (vent à 10~m), \texttt{psl} (pression au niveau de la mer) et \texttt{huss} (humidité spécifique). Un champ est représenté comme un tableau $[\text{canaux} \times \text{temps} \times \text{lat} \times \text{lon}]$.

\subsection{Propriétés statistiques de la précipitation}
\label{ssec:stats-precip}

La précipitation quotidienne présente trois propriétés qui contraignent toute méthode de downscaling. (i)~\textbf{Intermittence} : 50 à 80\,\% des jours selon les régions ne reçoivent aucune précipitation détectable, créant un pic de probabilité en zéro incompatible avec l'hypothèse de variable continue. (ii)~\textbf{Asymétrie à queue lourde} : conditionnellement à un jour pluvieux, l'intensité suit approximativement une loi log-normale, si bien que toute méthode gaussienne sous-estime systématiquement les extrêmes~\cite{maraun_widmann_2018}, précisément les événements à fort enjeu. (iii)~\textbf{Hétérogénéité spatiale} : la précipitation est structurée par le relief, la ligne de côte et l'occupation du sol, avec des gradients de plusieurs cm/km en zone de relief --- structure invisible à basse résolution et argument central en faveur du downscaling. Toute méthode pertinente doit reproduire ces trois traits ; le pipeline expérimental, détaillé au chapitre Matériels et méthodes, applique un régridage LR ($23 \times 26$ pixels sur la Nouvelle-Zélande) et fournit les cibles HR ($172 \times 179$ pixels).
\label{fig:lognormal}

% =============================================================
\section{Approches classiques de downscaling}
\label{sec:classiques}
% =============================================================

Le downscaling comporte une dimension \emph{spatiale} (grille grossière $\to$ grille fine) et une dimension \emph{temporelle} (fréquence faible $\to$ fréquence élevée). Notre contribution couvre exclusivement la dimension spatiale ; la dimension temporelle est brièvement évoquée pour positionner le champ.

\subsection{Désagrégation spatiale}
\label{ssec:disagg-spatiale}

\subsubsection{Statistique classique}

Les approches statistiques apprennent une relation prédicteurs LR $\to$ prédictand HR : \textbf{régression linéaire multiple} pixel-par-pixel, \emph{Bias Correction} (BC) et \emph{Quantile Mapping} (QM) alignant les distributions marginales sur les observations~\cite{wilby_wigley_1997,maraun_widmann_2018}. Simples, interprétables, peu coûteuses, ces méthodes butent sur trois limites structurelles : hypothèse de \emph{stationnarité} contestable sous changement climatique, perte de la cohérence spatiale (traitement pixel-par-pixel), absence native d'incertitude calibrée.

\subsubsection{Dynamique (RCM)}

L'approche dynamique exécute un RCM emboîté dans un GCM (conditions aux limites latérales) qui résout les équations physiques à $10$--$50\,\mathrm{km}$ sur un domaine restreint. Elle bénéficie de la cohérence physique et d'une transférabilité de principe hors-distribution, mais hérite des biais du GCM forçant, mobilise des ressources de calcul considérables et reste hors de portée des équipes du Sud.

\subsubsection{Apprentissage supervisé déterministe}

L'apprentissage profond a renouvelé l'approche statistique : \textbf{forêts aléatoires} et \texttt{XGBoost} pour la prédiction point-par-point, puis CNN de super-résolution (SRCNN, U-Net, EDSR) traitant le champ HR comme un tout cohérent. Ces modèles restent fondamentalement \emph{déterministes} : une sortie unique par entrée. Or le downscaling est intrinsèquement mal posé --- de nombreux champs HR sont compatibles avec une même entrée LR ---, si bien que la sortie déterministe se réduit à une moyenne conditionnelle qui lisse les extrêmes et sous-estime la variabilité. C'est ce constat qui motive le glissement vers les approches probabilistes (section~\ref{sec:generatifs}).

\subsection{Désagrégation temporelle}
\label{ssec:disagg-temporelle}

Mathématiquement analogue à la désagrégation spatiale, la désagrégation temporelle (fréquence mensuelle $\to$ quotidienne ou horaire) pose des enjeux propres : autocorrélation, conservation des cumuls, extrêmes instantanés. Les approches historiques (cascades multiplicatives, fragmentation, chaînes de Markov) et profondes (RNN, LSTM, Transformers temporels) restent pour l'essentiel déterministes et ponctuelles, avec peu de cadres unifiés spatio-temporels. Le présent travail couvre exclusivement la dimension spatiale ; l'extension temporelle (module récurrent ou attentif) est laissée aux travaux futurs.

% =============================================================
\section{Approches probabilistes et génératives}
\label{sec:generatifs}
% =============================================================

\subsection{Du déterministe au probabiliste}
\label{ssec:det-vs-proba}

L'insuffisance des méthodes déterministes (section~\ref{ssec:disagg-spatiale}) n'est pas une difficulté technique mais une conséquence du caractère multi-valué du problème : à une carte LR correspond une famille de cartes HR plausibles, dont une seule a été observée. La reformulation probabiliste vise donc la \emph{distribution conditionnelle}
\begin{equation}
  p\!\left(\xHR \,\big|\, \yLR,\, s_{\HR}\right),
  \label{eq:objectif-proba}
\end{equation}
où $s_{\HR}$ regroupe d'éventuels champs auxiliaires HR (orographie, occupation du sol). On peut alors produire des \emph{ensembles} d'échantillons, calculer des quantiles et évaluer la couverture des extrêmes. La décomposition propre à \Oracle{} (moyenne causale + résiduel stochastique) est détaillée au chapitre dédié.

\subsection{Modèles génératifs adversariaux}
\label{ssec:gan}

Les premiers génératifs profonds appliqués au downscaling ont mobilisé les \textbf{GAN}~: générateur transformant un bruit en échantillon plausible, discriminateur distinguant générés et réels. Leur variante conditionnelle (cGAN) prend la carte LR comme conditionnement~; une référence néo-zélandaise pour la précipitation est~\cite{rampal_cgan_2024}, qui combine un générateur résiduel et une perte adversariale pondérée. Trois limites empiriques sont bien documentées : (i)~\textbf{sub-dispersion} (variance prédite $<$ observée, sous-estimation des extrêmes), (ii)~\textbf{forte sensibilité} au coefficient adversarial $\lambda_{\mathrm{adv}}$, (iii)~risque de \textbf{mode collapse}. Ces limites ont motivé l'exploration des modèles de diffusion.
\label{fig:subdispersion}

\subsection{Modèles de diffusion}
\label{ssec:diffusion}

Un modèle de diffusion définit un processus direct qui bruite progressivement le signal, puis apprend à l'inverser. Cette formulation, due aux DDPM de Ho et~al.~\cite{ho_ddpm_2020}, a été généralisée par Karras et~al. sous le nom d'EDM~\cite{karras_edm_2022}, qui introduit un préconditionnement systématique pour un entraînement stable et un échantillonnage de qualité en peu d'étapes. Pour le downscaling, la diffusion apporte trois avantages sur les GAN : entraînement stable (pas de déséquilibre adversarial), distribution de sortie expressive (modes multiples), meilleure calibration probabiliste. Le formalisme EDM et l'échantillonnage de Heun sont détaillés au chapitre suivant.

\subsection{Le paradigme deux étages : CorrDiff}
\label{ssec:corrdiff}

\CorrDiff{}~\cite{mardani_corrdiff_2023}, proposé par NVIDIA Research, découple explicitement ce qui est facile à apprendre et ce qui est intrinsèquement stochastique : un premier étage régresse la moyenne HR attendue, un second échantillonne par diffusion le résiduel par rapport à cette moyenne. Le champ HR s'écrit comme la somme d'une moyenne déterministe et d'un résiduel stochastique (formalisation au chapitre suivant).

\begin{figure}[ht]
  \centering
  \figplaceholder{Figure F4 à produire --- schéma horizontal en deux étages : (i) entrée LR + sur-échantillonnage $\to$ réseau de régression $\to$ moyenne $\mu_{HR}$ ; (ii) bruit + conditionnement $\to$ réseau de diffusion $\to$ résiduel $\delta$ ; addition finale $\to$ sortie HR.}
  \caption{Paradigme deux étages de Mardani et al.~\cite{mardani_corrdiff_2023} : régression de $\muHR$ puis diffusion du résiduel $\delta$.}
  \label{fig:corrdiff}
\end{figure}

L'intérêt est triple : (i)~\emph{stabilité} (régression déterministe convergente, diffusion sur résidus quasi-gaussiens) ; (ii)~\emph{calibration} (variance explicitement contrôlée par le second étage, non plus produit d'un équilibre adversarial) ; (iii)~\emph{efficacité} (diffusion limitée au résiduel de faible amplitude). Ce paradigme constitue le point de départ direct de notre contribution. \Oracle{} en reprend la structure mais introduit une contrainte causale au premier étage : $\muHR$ y est prédite à travers un graphe causal appris, ce qui en fait par construction une version \emph{causalisée} de \CorrDiff{}.

\paragraph{Évaluation probabiliste.}
\label{ssec:eval-encart}
Les métriques déterministes classiques (RMSE, MAE, Pearson) ne suffisent pas à évaluer un modèle probabiliste, dont la valeur réside dans la distribution prédite. La littérature mobilise un panel élargi : \emph{CRPS} (généralisation probabiliste de l'erreur absolue), \emph{spread/skill} (calibration du second moment), \emph{RAPSD} (fidélité du spectre spatial), \emph{LHD} (distributions locales). La présentation formelle et l'application à \Oracle{} sont reportées au chapitre Matériels et méthodes.

% =============================================================
\section{Causalité et modèles structurels}
\label{sec:causalite}
% =============================================================

\subsection{Pourquoi la causalité en downscaling climatique}
\label{ssec:why-causal}

Les approches probabilistes précédentes sont \emph{corrélatives} : elles apprennent une distribution conditionnelle reproduisant les co-occurrences observées, ce qui ne garantit pas la qualité prédictive hors du régime d'entraînement --- précisément la situation du changement climatique. La causalité vise à dépasser cette limite par trois apports : (i)~\textbf{robustesse OOD}, un modèle reposant sur des relations causales identifiées (e.g.~influence de l'humidité basse couche sur la convection) étant en principe plus stable sous changement de régime ; (ii)~\textbf{interprétabilité physique}, la structure causale étant confrontable à la connaissance experte ; (iii)~\textbf{capacité d'intervention}, le modèle répondant à des questions \og que se passerait-il si \dots \fg{} sans recours à des analogues.

\subsection{Modèles causaux structurels et graphes acycliques dirigés}
\label{ssec:scm-dag}

Le cadre formel de la causalité est celui des \textbf{modèles causaux structurels} (SCM) de Pearl~\cite{pearl_causality_2009,pearl_book_of_why_2018}. Un SCM est un triplet $(V, U, F)$ où $V$ regroupe les variables observables, $U$ des bruits exogènes indépendants et $F$ des fonctions structurelles telles que
\begin{equation}
  V_i \;=\; f_i\!\left(\mathrm{Pa}(V_i),\, U_i\right),
  \label{eq:scm}
\end{equation}
$\mathrm{Pa}(V_i)$ étant l'ensemble des parents causaux de $V_i$. Le SCM s'exprime naturellement par un \emph{Directed Acyclic Graph} (DAG) : nœuds = variables, arcs = liens parent~$\to$~enfant. L'acyclicité garantit l'identifiabilité (ordre topologique d'évaluation) et fournit une description visuelle des dépendances --- par exemple, sur un mini-système atmosphérique, humidité, vent, température et orographie déterminent la précipitation locale, sans rétroaction explicite.
\label{fig:scm-dag}

\subsection{Intervention, identifiabilité et lien avec le downscaling}
\label{ssec:intervention}

L'opérateur $\mathrm{do}$ de Pearl définit l'\textbf{intervention} : $\mathrm{do}(X = x)$ fixe la valeur de $X$ indépendamment de ses parents (suppression de ses flèches entrantes) et fournit une distribution $p(Y\,|\,\mathrm{do}(X = x))$ en général distincte de la conditionnelle naïve $p(Y\,|\,X = x)$. C'est cette distinction qui sépare un modèle corrélatif d'un modèle structurellement causal. Un modèle \emph{intervention-aware} a un effet calculable et prévisible sous intervention sur ses entrées. Cette propriété (objectif O3 d'\Oracle{}, chapitre dédié) est obtenue \emph{au sens structurel} via la matrice $A_{\mathrm{dag}}$ apprise. Elle ne se confond pas avec une causalité \emph{physique} (où les flèches coïncideraient avec les lois atmosphériques fait pour fait).

\label{ssec:learned-dag}
Apprendre la structure d'un DAG depuis les données est un problème combinatoire dur. Bello, Aragam et Ravikumar~\cite{bello_dagma_2022} ont proposé \emph{DAGMA}, une caractérisation continue et différentiable de l'acyclicité (déterminant logarithmique d'une M-matrice), intégrable directement dans l'entraînement d'un réseau profond. \Oracle{} mobilise DAGMA pour insérer la causalité structurelle dans un paradigme génératif à deux étages style \CorrDiff{}, en prédisant $\muHR$ comme produit causal du DAG appris ; détail formel et examen comparatif sont reportés au chapitre dédié.

% =============================================================
\section{Verrous scientifiques et trilemme}
\label{sec:verrous}
% =============================================================

Les fondements précédents cristallisent trois verrous structurants du downscaling génératif. Nous les présentons avant de les articuler dans un \emph{trilemme}.

\subsection{Verrou 1 : stabilité d'entraînement}
\label{ssec:verrou-stabilite}

Les modèles adversariaux sont très sensibles au coefficient $\lambda_{\mathrm{adv}}$ et au déséquilibre générateur-discriminateur, difficiles à régler de manière transférable d'une région à l'autre~\cite{rampal_cgan_2024}. Les modèles de diffusion, plus stables en principe, exigent néanmoins une calibration soignée du préconditionnement et du schéma de bruit --- c'est précisément ce qu'apporte EDM~\cite{karras_edm_2022}. Ce verrou conditionne reproductibilité et déployabilité multi-cas.

\subsection{Verrou 2 : réalisme des extrêmes et sub-dispersion}
\label{ssec:verrou-extremes}

Le verrou le plus critique pour nos applications est le \textbf{réalisme des extrêmes}, via le phénomène de \emph{sub-dispersion} : variance prédite conditionnellement à l'entrée inférieure à la variance observée. Mesurée par le rapport $\sigma_{\mathrm{pred}} / \mathrm{RMSE}$, elle vaut $\approx 1$ pour un modèle bien calibré, mais des valeurs autour de $0{,}3$ sont rapportées en downscaling génératif~\cite{rampal_cgan_2024,mardani_corrdiff_2023,harris_2022}, soit une sous-estimation d'un facteur trois de la variabilité fine, donc des extrêmes --- précisément les événements à plus forte valeur ajoutée. Sa mesure sur \Oracle{} et ses ablations causales sera l'un des objets centraux de l'évaluation.

\subsection{Verrou 3 : robustesse hors-distribution}
\label{ssec:verrou-ood}

Sous changement climatique, la statistique conjointe des variables atmosphériques évolue : un modèle historique doit conserver sa qualité hors des analogues observés. Les modèles purement corrélatifs y sont vulnérables (associations contingentes au régime d'entraînement), tandis que la causalité structurelle (section~\ref{sec:causalite}) apparaît comme une réponse de principe : un modèle organisé autour de relations causales doit, en théorie, mieux transférer en \textbf{OOD}.

\subsection{Trilemme et positionnement d'\Oracle{}}
\label{ssec:trilemme}

L'examen de la littérature récente montre que les approches qui réussissent sur l'un de ces axes échouent souvent sur au moins un autre : les GAN les plus expressifs sont les plus instables, les modèles de diffusion stables restent corrélatifs et sans garantie OOD, les RCM dynamiques robustes sont coûteux et peu transférables. Cette tension, présente sous diverses formes chez plusieurs auteurs~\cite{rampal_cgan_2024,mardani_corrdiff_2023,harris_2022}, peut être condensée en un \emph{trilemme} entre stabilité d'entraînement, réalisme des extrêmes et robustesse OOD (figure~\ref{fig:trilemme}). Nous le présentons \emph{tel qu'il émerge de la littérature}, et non comme une formalisation propre.

\begin{figure}[ht]
  \centering
  \begin{tikzpicture}[
    vertex/.style={font=\small\bfseries, text=uridarkblue, align=center},
    method/.style={circle, draw, thick, inner sep=2.5pt, font=\footnotesize\sffamily},
    methodlbl/.style={font=\scriptsize\sffamily, align=center},
  ]
    % Triangle (equilateral, side ~6cm)
    \coordinate (top)   at (0, 3.46);   % stabilité (haut)
    \coordinate (left)  at (-3, -1.5);  % réalisme extrêmes (bas-gauche)
    \coordinate (right) at (3, -1.5);   % robustesse OOD (bas-droite)

    \draw[thick, uridarkblue, rounded corners=1pt]
      (top) -- (left) -- (right) -- cycle;

    % Etiquettes sommets
    \node[vertex, above=0.15cm of top] {Stabilité \\ d'entraînement};
    \node[vertex, below left=0.05cm and -0.4cm of left] {Réalisme \\ des extrêmes};
    \node[vertex, below right=0.05cm and -0.4cm of right] {Robustesse \\ OOD};

    % Points : approches positionnées qualitativement à l'intérieur
    %  - cGAN : bon en extrêmes, faible en stabilité, faible en OOD
    %  - CorrDiff : bon en stabilité, bon en extrêmes, faible OOD
    %  - Oracle : bonne stabilité + extrêmes corrects + gain OOD
    %  - Idéal : centroïde (équilibre des 3)

    \node[method, fill=red!25] (cgan) at (-1.6, -0.6) {};
    \node[methodlbl, below=2pt of cgan, text=red!60!black] {cGAN};

    \node[method, fill=orange!40] (corrdiff) at (0.7, 1.0) {};
    \node[methodlbl, above right=-2pt and 0pt of corrdiff, text=orange!70!black] {CorrDiff};

    \node[method, fill=uriaccent!70] (oracle) at (0.4, 0.05) {};
    \node[methodlbl, right=2pt of oracle, text=uriaccent] {\Oracle{}};

    \node[method, fill=green!55!black, text=white] (ideal) at (0, 0.15) {};
    \node[methodlbl, below=2pt of ideal, text=green!45!black] {Idéal};
  \end{tikzpicture}
  \caption{Trilemme du downscaling génératif. Les trois sommets représentent les verrous identifiés ; les points positionnent qualitativement les principales approches. Plus un point est proche d'un sommet, mieux l'approche traite ce verrou. Constat agrégé à partir de~\cite{rampal_cgan_2024,mardani_corrdiff_2023,harris_2022}.}
  \label{fig:trilemme}
\end{figure}
\label{ssec:positionnement-oracle}

\Oracle{} se positionne dans ce trilemme en s'inscrivant dans le paradigme à deux étages de \CorrDiff{} (stabilité, verrou~1), en y introduisant une contrainte structurellement causale au premier étage (robustesse OOD, verrou~3), et en mobilisant la diffusion EDM au second étage pour la calibration probabiliste (réalisme des extrêmes, verrou~2). La figure~\ref{fig:cartouche-oracle} en propose une vue d'ensemble ; détails architecturaux et formalisation des objectifs sont déférés au chapitre dédié.

\begin{figure}[ht]
  \centering
  \resizebox{\linewidth}{!}{%
  \begin{tikzpicture}[
    block/.style={rectangle, draw=uridarkblue, fill=uridarkblue!8, rounded corners=2pt,
                  minimum height=1.1cm, minimum width=2.1cm, align=center, font=\small,
                  inner sep=3pt},
    cblock/.style={rectangle, draw=uriaccent, fill=uriaccent!10, rounded corners=2pt,
                   minimum height=1.1cm, minimum width=2.1cm, align=center, font=\small,
                   inner sep=3pt, very thick},
    outblock/.style={rectangle, draw=urigray, fill=urigray!10, rounded corners=2pt,
                     minimum height=1.0cm, minimum width=5.0cm, align=center, font=\small},
    stagelbl/.style={font=\footnotesize\bfseries\itshape, text=urigray, align=center},
    arr/.style={-{Stealth[length=2.2mm]}, thick, uridarkblue!75},
    carr/.style={-{Stealth[length=2.2mm]}, thick, uriaccent!90},
    causal/.style={font=\footnotesize\bfseries, text=uriaccent},
  ]
    % --- Stage 1 row (causalisé) ---
    \node[block]                       (lr)  at (0,0)          {$\LR$\\ \scriptsize ACCESS-CM2};
    \node[cblock, right=0.7cm of lr]   (enc)                   {Encodeur\\ \scriptsize intelligible};
    \node[cblock, right=0.7cm of enc]  (dag)                   {Graphe causal\\ $\Adag$ \scriptsize (DAGMA)};
    \node[cblock, right=0.7cm of dag]  (mu)                    {$\muHR$\\ \scriptsize moyenne};

    % --- Stage 2 row (diffusion EDM) ---
    \node[block, below=2.4cm of enc]   (noise)                 {Bruit $+$\\ $\muHR$};
    \node[block, right=0.7cm of noise] (edm)                   {Diffusion\\ EDM};
    \node[block, right=0.7cm of edm]   (delta)                 {Résidu $\widehat{\delta}$};

    % --- Output row ---
    \node[outblock, below=1.3cm of edm] (out)
      {$\widehat{x}_{\HR} \,=\, \mathrm{baseline} \,+\, \muHR \,+\, \widehat{\delta}$};

    % --- Stage frames + labels ---
    \node[draw=uriaccent, dashed, rounded corners=4pt, inner sep=7pt,
          fit=(lr) (enc) (dag) (mu)] (s1box) {};
    \node[stagelbl, text=uriaccent, anchor=south west]
      at (s1box.north west) {Stage 1 — causalisation \Oracle{}};

    \node[draw=urigray, dashed, rounded corners=4pt, inner sep=7pt,
          fit=(noise) (edm) (delta)] (s2box) {};
    \node[stagelbl, anchor=south west]
      at (s2box.north west) {Stage 2 — diffusion résiduelle EDM (paradigme \CorrDiff{})};

    % --- Arrows Stage 1 ---
    \draw[carr] (lr)  -- (enc);
    \draw[carr] (enc) -- (dag);
    \draw[carr] (dag) -- (mu);

    % --- mu_HR -> noise (passe vers stage 2) ---
    \draw[arr, dashed] (mu.south) -- ++(0, -0.6) -| (noise.north);

    % --- Stage 2 row arrows ---
    \draw[arr] (noise) -- (edm);
    \draw[arr] (edm)   -- (delta);

    % --- Convergence vers la sortie HR ---
    \draw[arr] (delta.south) -- ++(0, -0.5) -| (out.north east);
    \draw[arr] (mu.south)    -- ++(0, -0.3) -| ([xshift=2.0cm]out.north);
    \draw[arr] (lr.west)     -- ++(-0.6, 0) |- (out.west)
       node[pos=0.75, above, font=\scriptsize, text=urigray] {baseline};

    % --- Causalisation annotation ---
    \node[causal, above=0.45cm of dag] {causalisation};
    \draw[uriaccent, thick, ->] ([yshift=2pt]dag.north) -- ++(0, 0.30);
  \end{tikzpicture}%
  }
  \caption{Cartouche conceptuel d'\Oracle{} : version causalisée du paradigme deux étages de \CorrDiff{}. L'étage 1 estime $\muHR$ via un graphe causal appris ; l'étage 2 échantillonne par diffusion le résiduel $\delta$.}
  \label{fig:cartouche-oracle}
\end{figure}

% =============================================================
\section{Synthèse}
\label{sec:synthese}
% =============================================================

Ce chapitre a établi trois acquis. D'abord, le besoin de désagrégation est concret et localisé : les GCM fournissent une information à une résolution incompatible avec les usages locaux, particulièrement en contexte africain où l'accès aux RCM dynamiques reste limité ; la précipitation y est la cible à la fois la plus pertinente et la plus exigeante. Ensuite, le paradigme méthodologique de référence a glissé du déterministe vers le probabiliste, puis le génératif, et plus récemment vers le paradigme à deux étages de \CorrDiff{} (moyenne déterministe + résiduel par diffusion), qui constitue le point de départ direct de notre travail. Enfin, trois verrous structurants émergent --- stabilité d'entraînement, réalisme des extrêmes (sub-dispersion), robustesse OOD --- articulés dans un trilemme auquel \Oracle{} répond en greffant une structure causale sur le paradigme à deux étages.

Le chapitre suivant examine l'état de l'art : travaux antérieurs dans la lignée du downscaling probabiliste à deux étages (\CorrDiff{}, cGAN de Rampal et~al.) et approches éclairant l'intégration de la causalité au cadre génératif.

% =============================================================
% Bibliographie (placeholder, à migrer vers un .bib en phase finale)
% =============================================================
\begin{thebibliography}{99}

\bibitem{ipcc_ar6_2021} IPCC. \emph{Climate Change 2021: The Physical Science Basis. Contribution of Working Group~I to the Sixth Assessment Report of the Intergovernmental Panel on Climate Change}. Cambridge University Press, 2021.

\bibitem{almazroui_2020} M.~Almazroui, S.~Saeed, F.~Saeed, M.~N.~Islam, and M.~Ismail. Projections of precipitation and temperature over the South Asian countries in CMIP6. \emph{Earth Systems and Environment}, 4:297--320, 2020.

\bibitem{eyring_cmip6_2016} V.~Eyring, S.~Bony, G.~A.~Meehl, C.~A.~Senior, B.~Stevens, R.~J.~Stouffer, and K.~E.~Taylor. Overview of the Coupled Model Intercomparison Project Phase 6 (CMIP6) experimental design and organization. \emph{Geoscientific Model Development}, 9:1937--1958, 2016.

\bibitem{hersbach_era5_2020} H.~Hersbach et al. The ERA5 global reanalysis. \emph{Quarterly Journal of the Royal Meteorological Society}, 146:1999--2049, 2020.

\bibitem{maraun_widmann_2018} D.~Maraun and M.~Widmann. \emph{Statistical Downscaling and Bias Correction for Climate Research}. Cambridge University Press, 2018.

\bibitem{vandal_deepsd_2017} T.~Vandal, E.~Kodra, S.~Ganguly, A.~Michaelis, R.~Nemani, and A.~R.~Ganguly. DeepSD: Generating high resolution climate change projections through single image super-resolution. In \emph{Proceedings of the 23rd ACM SIGKDD International Conference on Knowledge Discovery and Data Mining}, 2017.

\bibitem{wmo_gcos_2022} World Meteorological Organization. \emph{The 2022 GCOS Implementation Plan} (GCOS-244). World Meteorological Organization, Geneva, 2022.

\bibitem{wilby_wigley_1997} R.~L.~Wilby and T.~M.~L.~Wigley. Downscaling general circulation model output: a review of methods and limitations. \emph{Progress in Physical Geography}, 21(4):530--548, 1997.

\bibitem{ho_ddpm_2020} J.~Ho, A.~Jain, and P.~Abbeel. Denoising diffusion probabilistic models. In \emph{Advances in Neural Information Processing Systems}, 2020.

\bibitem{karras_edm_2022} T.~Karras, M.~Aittala, T.~Aila, and S.~Laine. Elucidating the design space of diffusion-based generative models. In \emph{Advances in Neural Information Processing Systems}, 2022.

\bibitem{mardani_corrdiff_2023} M.~Mardani, N.~Brenowitz, Y.~Cohen, J.~Pathak, C.-Y.~Chen, C.-C.~Liu, A.~Vahdat, K.~Kashinath, J.~Kautz, and M.~Pritchard. Residual corrective diffusion modeling for km-scale atmospheric downscaling. \emph{arXiv preprint arXiv:2309.15214}, 2023.

\bibitem{rampal_cgan_2024} N.~Rampal, P.~B.~Gibson, S.~Hobeichi, J.~Sherwood, and S.~B.~Power. Reliable generative downscaling of climate data with conditional GANs. \emph{Journal of Advances in Modeling Earth Systems}, 2024.

\bibitem{harris_2022} L.~Harris, A.~T.~T.~McRae, M.~Chantry, P.~D.~Dueben, and T.~N.~Palmer. A generative deep learning approach to stochastic downscaling of precipitation forecasts. \emph{Journal of Advances in Modeling Earth Systems}, 14(10):e2022MS003120, 2022.

\bibitem{pearl_causality_2009} J.~Pearl. \emph{Causality: Models, Reasoning, and Inference}. Cambridge University Press, 2nd edition, 2009.

\bibitem{pearl_book_of_why_2018} J.~Pearl and D.~Mackenzie. \emph{The Book of Why: The New Science of Cause and Effect}. Basic Books, 2018.

\bibitem{bello_dagma_2022} K.~Bello, B.~Aragam, and P.~Ravikumar. DAGMA: Learning DAGs via M-matrices and a log-determinant acyclicity characterization. In \emph{Advances in Neural Information Processing Systems}, 2022.

\end{thebibliography}

\end{document}
